\section{Optimisation: SGD, the schedule and the training loop}
\label{sec:optim}

\subsection{The optimiser}
\label{sec:sgd}

\begin{repo}[\texttt{scripts/train\_resnet.py}]
\begin{lstlisting}[firstnumber=56]
optim = torch.optim.SGD(params=model.parameters(),lr=0.1,momentum=0.9,weight_decay=1e-4)
loss_fct = nn.CrossEntropyLoss()
\end{lstlisting}
One parameter group containing all 269{,}722 parameters, with \code{dampening=0} and \code{nesterov=False} (defaults).
\end{repo}

\begin{idea}[the exact update PyTorch performs]
At step $t$, for every parameter tensor $\theta$ (all shapes are the parameter's own shape):
\begin{align*}
  g_t &= \nabla_\theta\mathcal L(\theta_t) + \lambda\,\theta_t && \text{weight decay folded into the gradient}\\
  v_t &= \begin{cases} g_t & t=1\\ \mu\,v_{t-1} + g_t & t>1\end{cases} && \text{momentum buffer (velocity)}\\
  \theta_{t+1} &= \theta_t - \eta_t\,v_t && \text{step}
\end{align*}
with $\mu=0.9$, $\lambda=10^{-4}$ and $\eta_t$ from the schedule. The momentum buffer is initialised to the first gradient, not to zero. Unrolled, $v_t=\sum_{k\le t}\mu^{t-k}g_k$ is an exponentially weighted sum of past gradients with an effective memory of $1/(1-\mu)=10$ steps. Under a constant gradient $g$, $v\to g/(1-\mu)=10g$, so \textbf{the effective step size is $10\eta$}: $1.0$, then $0.1$, then $0.01$.
\end{idea}

\begin{warn}[two momentum conventions]
Some papers and frameworks (e.g.\ Sutskever et al.\ 2013, Caffe) put the learning rate \emph{inside} the velocity, $v_t=\mu v_{t-1}-\eta_t g_t$, $\theta_{t+1}=\theta_t+v_t$. With a constant $\eta$ the two are identical ($v^\text{Caffe}=-\eta\,v^\text{PyTorch}$). They differ only at a learning-rate drop. In PyTorch the stored velocity is unchanged and gets multiplied by the new, smaller $\eta$, so the step shrinks $10\times$ immediately. In the other form the old velocity still carries the old $\eta$ and fades over $\sim$20 steps. Over 16{,}000-iteration phases the difference is negligible, but it is worth knowing when reading other code.
\end{warn}

\begin{repo}[one step verified]
Three steps of the update above, written out by hand, reproduce \code{torch.optim.SGD(lr=0.1, momentum=0.9, weight\_decay=1e-4)} to within $4\times10^{-6}$ (\code{float32}). The optimiser state is one momentum buffer per parameter tensor, with the same shape, i.e.\ another 269{,}722 numbers. That is why the training checkpoints (2.2\,MB) are twice the size of the exported serving weights (1.1\,MB).
\end{repo}

\subsection{Weight decay}
\label{sec:wd}

\begin{idea}[L2 regularisation, and what it does with BN]
Adding $\lambda\theta$ to the gradient is gradient descent on
\[
  \mathcal L(\theta) + \tfrac{\lambda}{2}\|\theta\|_2^2 ,
\]
so with plain SGD, ``weight decay'' and ``L2 penalty'' are the same thing. (They differ only for adaptive optimisers like Adam, hence AdamW.) Each step shrinks every weight by the factor $(1-\eta\lambda)$ before the gradient step: $1-10^{-5}$ per step at $\eta=0.1$.

In a network where every conv is followed by BN, the conv's output is invariant to rescaling its weights ($\BN(\alpha\,\mathcal W*x)=\BN(\mathcal W*x)$ for $\alpha>0$). The penalty therefore does not change the function, but it keeps $\|\mathcal W\|$ small, and a small $\|\mathcal W\|$ makes the same $\eta$ a \emph{larger relative} step. Weight decay on BN-ed convs acts largely as a control on the effective learning rate.
\end{idea}

\begin{repo}[applied to every parameter]
Because \code{model.parameters()} goes into a single group, decay applies to conv weights, BN $\gamma$ \emph{and} $\beta$, and the linear weight and bias alike. The paper does not say whether BN parameters were excluded; many reimplementations decay everything, as here. Measured on the final model: $\|\theta\|_2^2 = 3218.6$, so the penalty term is $\tfrac{\lambda}{2}\|\theta\|^2 = 0.161$. At the end of training the data loss is $0.024$, so \textbf{the regulariser is about $6.6\times$ the data term} of the objective SGD is actually minimising. Its visible fingerprint is stem channel 0 (\cref{sec:forward}): its filter norm decayed to $0.047$ and its $\gamma$ to $0.011$, a channel the data did not defend.
\end{repo}

\subsection{The learning-rate schedule}
\label{sec:schedule}

\begin{lstlisting}[firstnumber=60]
def lr_at(iteration):
    if iteration < 32000:
        return 0.1
    elif iteration < 48000:
        return 0.01
    else:
        return 0.001
def set_lr(optimizer,lr):
    for p in optimizer.param_groups:
        p["lr"] = lr
\end{lstlisting}

\begin{idea}[step decay, counted in iterations]
$\eta_t = 0.1\cdot 10^{-\,[t\ge32000]-[t\ge48000]}$, where $t$ is the number of optimiser steps already taken. The paper specifies the schedule in iterations (``divide it by 10 at 32k and 48k iterations, and terminate training at 64k iterations''), independent of the dataset size, which is why the repo counts iterations, not epochs.
\end{idea}

\begin{center}
\begin{tikzpicture}
\begin{axis}[width=0.92\linewidth,height=4.2cm,ymode=log,xmin=0,xmax=64000,ymin=0.0005,ymax=0.2,
  xlabel={iteration $t$},ylabel={$\eta_t$},xtick={0,16000,32000,48000,64000},
  xticklabels={0,16k,32k,48k,64k},ytick={0.001,0.01,0.1},yticklabels={0.001,0.01,0.1},
  grid=major,grid style={black!10},tick label style={font=\scriptsize},label style={font=\small},scaled x ticks=false,
  extra x ticks={32000,48000},extra x tick labels={},extra x tick style={grid style={warn,dashed}}]
  \addplot[const plot,very thick,idea] coordinates {(0,0.1) (32000,0.01) (48000,0.001) (64000,0.001)};
  \node[font=\scriptsize,align=center] at (axis cs:16000,0.03) {32{,}000 steps at 0.1\\(epochs 1--82 on 50k)};
  \node[font=\scriptsize,align=center] at (axis cs:40000,0.0032) {16{,}000 at 0.01\\(epochs 82--123)};
  \node[font=\scriptsize,align=center] at (axis cs:56000,0.0032) {16{,}000 at 0.001\\(epochs 123--164)};
\end{axis}
\end{tikzpicture}
\end{center}

\begin{repo}[iterations $\leftrightarrow$ epochs]
\begin{center}
\small
\begin{tabular}{@{}lrrrrr@{}}
\toprule
run & iters/epoch & 1st drop (32k) & 2nd drop (48k) & end (64k) & last epoch\\
\midrule
final (50k) & 391 & during epoch 82 & during epoch 123 & epoch 164 & 267 iterations\\
development (45k) & 352 & during epoch 91 & during epoch 137 & epoch 182 & 288 iterations\\
\bottomrule
\end{tabular}
\end{center}
$64{,}000\times128 = 8{,}192{,}000$ examples are seen, $163.84$ passes over 50k (counting the short last batch of each epoch as a full step, which is what the loop does). $163\cdot391=63{,}733$ iterations finish 163 epochs, and epoch 164 stops after $267$ more, exactly as in \file{artifacts/logs/final\_res20.csv} (last row: epoch 164, iteration 64{,}000). The learning rate is set \emph{before} every step from the current count, so step 0 through step 31{,}999 use 0.1 (\file{check\_schedule.py} tests the boundaries 31999/32000 and 47999/48000). Every BN layer's \code{num\_batches\_tracked} in the checkpoint reads exactly \textbf{64{,}000}: one train-mode forward per optimiser step, with no extra ones.
\end{repo}

\subsection{The training loop}
\label{sec:loop}

\begin{lstlisting}[firstnumber=33]
def train_one_epoch(model,loader,optimizer,loss_fn,device,iteration,max_iterations):
    model.train()
    num_exmpl  = 0
    correct_pred = 0
    loss_cum = 0
    iterer = iteration
    for image,label in loader:
        image = image.to(device)
        label = label.to(device)
        if iterer >= max_iterations:
             break
        set_lr(optimizer, lr_at(iterer))
        logits,loss = train_step(model=model,x=image,y=label,optimizer=optimizer,loss_fn=loss_fn)
        compare = (torch.argmax(logits,dim=1) == label)
        compare = compare.sum().item()
        correct_pred+= compare
        iterer+=1
        num_exmpl += logits.shape[0]
        loss_cum+=loss*logits.shape[0]
    loss_cum /= num_exmpl
    correct_pred /= num_exmpl
    return loss_cum,correct_pred,iterer
\end{lstlisting}

\begin{center}
\resizebox{\linewidth}{!}{\begin{tikzpicture}[
  n/.style={draw=ink,rounded corners=2pt,fill=#1,font=\footnotesize,minimum height=1.3cm,align=center,inner sep=3pt,text width=2.6cm},
  n/.default=white, ar/.style={-{Stealth[length=2mm]},thick}]
  \node[n=cIn] (a) at (0,0) {next batch\\$X$ \shp{128,3,32,32}\\$y$ \shp{128}};
  \node[n=cIn] (b) at (3.4,0) {\code{.to("cuda")}\\stop if $t\ge64000$};
  \node[n=cK] (c) at (6.8,0) {$\eta\leftarrow$ \code{lr\_at}$(t)$};
  \node[n=c32] (d) at (10.2,0) {\code{zero\_grad}\\forward, $\mathcal L$\\backward, step};
  \node[n=cG] (e) at (13.6,0) {sum \#correct\\and $\mathcal L\cdot|B|$\\$t\leftarrow t+1$};
  \foreach \p/\q in {a/b,b/c,c/d,d/e} \draw[ar] (\p) -- (\q);
  \draw[ar] (e.south) -- ++(0,-0.45) -| (a.south);
  \node[font=\footnotesize,fill=white] at (6.8,-1.1) {391 times per epoch; then \code{evaluate()} on validation (dev runs) and one CSV row};
\end{tikzpicture}}
\end{center}

\begin{repo}[what the logged training metrics mean]
\begin{itemize}[nosep]
  \item \textbf{train\_loss}: $\sum_\text{batches}\mathcal L_B\,|B|\big/\sum|B|$, the exact per-example average over the epoch, correct also for the short last batch.
  \item \textbf{train\_acc}: computed from the logits of each step's forward pass, i.e.\ on \emph{augmented} images, with \emph{train-mode} BN, and \emph{before} that step's update. The model changes during the epoch, so this is a running average, not the accuracy of the end-of-epoch model. The README's ``train error'' column (e.g.\ 0.52\% for ResNet-20) is this quantity at the last epoch.
  \item \textbf{val\_loss / val\_acc}: from \code{evaluate()} after the epoch, in eval mode, on un-augmented validation images. Empty in the \code{--full} run, which has no validation set.
  \item \textbf{lr}: read from \code{optim.param\_groups[0]["lr"]} after the epoch, i.e.\ the rate of its last step.
\end{itemize}
\end{repo}

\begin{center}
\begin{tikzpicture}
\begin{axis}[width=0.95\linewidth,height=6.2cm,xmin=0,xmax=64000,ymin=0,ymax=30,
  xlabel={iteration},ylabel={error (\%)},xtick={0,16000,32000,48000,64000},xticklabels={0,16k,32k,48k,64k},scaled x ticks=false,
  grid=major,grid style={black!10},tick label style={font=\scriptsize},label style={font=\small},
  legend style={font=\scriptsize,at={(0.98,0.97)},anchor=north east,draw=black!20},legend cell align=left]
  \addplot[idea,thick] table[col sep=comma,x=iteration,y expr={100*(1-\thisrow{train_acc})}] {logs/final_res20.csv};
  \addlegendentry{final run (50k): train error (augmented, train mode)}
  \addplot[warn,thick] table[col sep=comma,x=iteration,y expr={100*(1-\thisrow{val_acc})}] {logs/dev_res20.csv};
  \addlegendentry{development run (45k): validation error}
  \addplot[warn!50,thin] table[col sep=comma,x=iteration,y expr={100*(1-\thisrow{train_acc})}] {logs/dev_res20.csv};
  \addlegendentry{development run (45k): train error}
  \draw[black!50,dashed] (axis cs:32000,0) -- (axis cs:32000,30);
  \draw[black!50,dashed] (axis cs:48000,0) -- (axis cs:48000,30);
\end{axis}
\end{tikzpicture}
\end{center}

The two learning-rate drops are the most visible events in training. At the first drop, the development run's validation error falls from $17.6\%$ (epoch 90, $\eta=0.1$, noisy) to $9.8\%$ one epoch later and $8.8\%$ the epoch after. The second drop gives a smaller, final improvement. At $\eta=0.1$ the validation curve is noisy because the effective step ($10\eta=1$) keeps the weights bouncing around a wide region. Smaller steps let SGD settle into it. The final run takes 2{,}182\,s ($\approx36$ minutes) on an RTX 3050, about 13\,s per epoch.

\subsection{Seeds, checkpoints and the run script}

\begin{repo}[sources of randomness, and what controls each]
\begin{center}
\small
\begin{tabularx}{\linewidth}{@{}>{\raggedright\arraybackslash}p{3.2cm}>{\raggedright\arraybackslash}X>{\raggedright\arraybackslash}X@{}}
\toprule
randomness & generator & seeded by\\
\midrule
weight initialisation & PyTorch global RNG & \code{torch.manual\_seed(args.seed)}, before the model is built\\
batch order & \code{torch.Generator g} passed to the loader & \code{g.manual\_seed(args.seed)}\\
crop offsets and flips & \code{np.random.default\_rng} in the dataset & \code{seed=args.seed}\\
train/val split & local \code{torch.Generator} & literal 0\\
GPU kernels & cuDNN algorithm choice & \emph{not controlled}\\
\bottomrule
\end{tabularx}
\end{center}
All reported runs use seed 0. Without \code{torch.backends.cudnn.deterministic=True}, re-running is reproducible up to floating-point noise, not bit-for-bit.
\end{repo}

\begin{repo}[checkpoints]
\code{save\_checkpoint} writes one dictionary,
\begin{center}\code{\{"model": state\_dict, "optimizer": state\_dict, "iteration": t, "seed": s\}},\end{center}
\emph{once}, at the end of training. Nothing is saved per epoch, so an interrupted run restarts from scratch. \code{load\_checkpoint} restores the model, restores the optimiser only if one is passed (the service passes \code{None}), and returns \code{(iteration, seed)}. \file{scripts/export\_model.py} re-saves the final weights without the optimiser state as \file{model/resnet20.pt} (1{,}121{,}617 bytes, 116 \code{state\_dict} entries) and copies the 50k mean next to it. Those two files are all the service needs.
\end{repo}

\begin{repo}[command line]
\code{python -m scripts.train\_resnet --name final\_res20 --full} for the headline model; \code{--n 9} for ResNet-56; \code{--plain} for the plain networks; \code{--max-iters} to shorten a run. The device is hard-coded to \code{"cuda"}.
\end{repo}

\subsection{Sanity checks that were run before the full training}

\begin{repo}[\texttt{scripts/}]
\begin{itemize}[nosep]
  \item \file{overfit\_batch.py}: 200 steps on one fixed batch of 128. A correct model, loss and optimiser must drive this to $\approx100\%$ training accuracy. If they cannot, something is broken before any hyperparameter matters. The first loss should be $\approx\ln10=2.30$. One random batch at initialisation gives $2.42$, slightly above because of the randomly initialised linear layer.
  \item \file{check\_model.py}: parameter counts for $n=3$ and $9$, residual and plain; output shape \shp{128,10}; initial weight standard deviations.
  \item \file{check\_schedule.py}: \code{lr\_at} on both sides of each boundary.
  \item \file{check\_checkpoint.py}: train 20 steps, save, load into a fresh model, and verify identical \code{state\_dict}s and identical validation metrics.
  \item \file{check\_eval.py}: train-mode forward passes followed by \code{evaluate()}, which must switch BN to eval mode.
  \item \file{check\_loader.py}, \file{check\_split.py}, \file{check\_mean.py}, \file{check\_augment.py}: data pipeline determinism and shapes, with the augmentation grid saved to \file{artifacts/aug\_grid.png}.
\end{itemize}
\end{repo}
